\section{Experimental Design}
\label{sec:design}

\subsection{Question and design}
\label{sec:design-question}

We ask whether building a human cognitive or developmental constraint into a language model is associated with closer correspondence to human behavioral and neural data than a standard, unconstrained LLM provides. We operationalize correspondence in two ways: predictive power, i.e.\ how well features derived from a frozen model predict held-out human responses, and representational similarity between model and response dissimilarity matrices. Neither is a test of shared mechanism.

Models are organized in a $2\times2$ design (Section~\ref{sec:design-models}): a memory factor ($M$; a memory-limited architecture or not) crossed with a data factor ($D$; a developmental, BabyLM-scale training corpus or web-scale text). This gives three questions: (H1) is a memory limitation associated with a change in alignment, (H2) is a developmental data budget, and (H3) do the two combine additively or interact. Because the four groups differ greatly in parameter count, depth and language-modeling quality (Table~\ref{tab:models}), a raw comparison of group means would confound the constraint with scale. We therefore compare groups conditional on size and perplexity (Section~\ref{sec:design-stats}), and answer each question separately for behavioral and neural responses.

Every dataset is crossed with every model and with two feature types (Section~\ref{sec:design-tracks}), giving a $2\times2$ of \{neural, behavioral\} data $\times$ \{embedding, surprisal\} predictors. The evaluation procedure is identical across models; only per-dataset configurations differ.

\subsection{Models}
\label{sec:design-models}

We evaluate 118 frozen models, organized into four groups (Table~\ref{tab:models}, Appendix~\ref{app:models}). Standard LLMs have $M=D=0$; Memory-Constrained and BabyLM models have $(M,D)=(1,0)$ and $(0,1)$; Constrained-BabyLM models have $M=D=1$. No model is fine-tuned on any behavioral or neural target.

\paragraph{Standard LLMs (59 models).} Off-the-shelf, web-scale autoregressive models with neither constraint: GPT-2, GPT-Neo, OPT, Pythia, OLMo-2, Qwen3, Llama-3, Ministral, and the recurrent and state-space models Mamba, Mamba2 and RWKV-4 (14M to 66B parameters, 6 to 65 layers), plus the two matched baselines released with the Forgetting Transformer (FoX), which have ordinary attention and no forget gate (760M parameters). Mamba, Mamba2 and RWKV-4 carry a fixed-size recurrent state but no constraint motivated by human memory; recurrent state alone does not define the tested memory-constraint category, so they are grouped here.

\paragraph{Memory-Constrained (24 models).} Architectures that limit how much of the preceding context the model can access or weight, motivated by capacity-limited accounts of human sentence processing such as activation-based retrieval \citep{lewisvasishth2005activation} and lossy-context surprisal \citep{futrell2020lossy}: distance-penalized attention (ALiBi), forget-gate and stick-breaking replacements for softmax attention, hard $n$-gram windows, a small bounded-state recurrent network (CBR-RNN), and web-scale models with the same mechanisms (BLOOM, whose attention uses ALiBi; FoX; MoLM). These are trained on web-scale text (15M to 7.1B parameters, 1 to 49 layers).

\paragraph{BabyLM (11 models).} Standard Transformer architectures trained on roughly the number of words a child hears by early adolescence, 10M to 100M words, following the BabyLM Challenge \citep{warstadt2023cfp,warstadt2023findings,hu2024secondbabylm} (33M to 140M parameters, 11 to 13 layers).

\paragraph{Constrained-BabyLM (24 models).} The ALiBi, forgetting, stick-breaking and hard $n$-gram architectures trained on the BabyLM corpora, i.e.\ both constraints at once (15M to 48M parameters, 2 to 4 layers).

Keeping the last group separate makes the crossing explicit and lets the regression ask whether the two constraints add up or interact. The categorization concerns the selected architectural mechanisms and training budgets, not a claim that all models within a group share one cognitive mechanism. Pythia and OLMo-2 intermediate training checkpoints are available but are not part of the comparison.

\subsection{Human data}
\label{sec:design-data}

Every dataset is reduced to the same unit of analysis, a word (or sentence, or item) embedded in its preceding context and paired with one or more measured human responses, so any model can be run through any dataset without dataset-specific model code. In total we use 18 datasets, which we split into 27 \emph{dataset units}: reading-time measures that come from the same recording (first-pass, go-past and scan-path times) and the three ZuCo reading tasks are counted as separate units, giving 16 behavioral and 11 neural units (Table~\ref{tab:datasets}, Appendix~\ref{app:datasets}). Units of one corpus share their items and participants and are therefore not independent; we return to this in Section~\ref{sec:design-stats}.

\paragraph{Behavioral datasets.}
Behavioral measures index processing difficulty indirectly, through the time a reader takes to move past a word.
\begin{itemize}
    \item \textbf{Self-paced reading.} Natural Stories \citep{futrell2021natural}, naturalistic narratives deliberately written to contain rare syntactic constructions, and UCL Reading \citep{frank2013reading}, 361 sentences read by 117 participants.
    \item \textbf{Eye-tracking.} GECO \citep{cop2017geco}, Provo \citep{luke2018provo}, the English part of MECO \citep{siegelman2022meco}, OneStop \citep{berzak2025onestop}, and the eye-tracking part of ZuCo \citep{hollenstein2018zuco}, where participants read isolated sentences under three task instructions (normal reading, sentiment judgment, task-specific relation judgment). Depending on the corpus we use first-pass duration, go-past time and scan-path (total dwell) time, each as its own unit.
\end{itemize}

\paragraph{Neural datasets.}
Neural measures record language processing from the brain at complementary spatial and temporal resolutions.
\begin{itemize}
    \item \textbf{fMRI.} Le Petit Prince \citep{li2022lepetitprince} (audiobook listening), Natural Stories fMRI \citep{shain2022robust} (responses pre-averaged within functionally localized language and multiple-demand regions), Tuckute et al.'s language-network dataset \citep{tuckute2024driving} (isolated sentences), and Pereira et al.\ \citep{pereira2018toward} (isolated sentences from two experiments).
    \item \textbf{EEG / ERP.} Broderick et al.'s continuous-listening audiobook corpus \citep{broderick2018electrophysiological}, and three N400 datasets built around cloze-probability manipulations: a nine-site replication of the classic article/noun expectancy design \citep{nieuwland2018largescale}, and two studies from the Szewczyk and Federmeier collection \citep{szewczyk2022context}, Federmeier et al.\ (2007) \citep{federmeier2007multiple} and Hubbard et al.\ (2019), analyzed separately.
    \item \textbf{Intracranial recordings.} A sentence-reading intracranial dataset \citep{fedorenko2016neural} and a large ECoG dataset recorded during a naturalistic podcast \citep{zada2025podcast}.
    \item \textbf{MEG.} MEG-MASC \citep{gwilliams2022neural}, passive listening to naturalistic stories.
\end{itemize}
Each unit retains its native word, sentence, item, or time-aligned response format.

\subsection{Final-layer predictors and direct prediction}
\label{sec:design-tracks}

For each model and dataset we compute two kinds of predictors from the same forward pass over the full preceding context. All reported analyses use the final hidden layer or the final output distribution. Models have between 1 and 65 layers, so ``the representation of a model'' needs a rule that is the same for all of them; the final layer involves no selection and treats shallow and deep models alike.

\paragraph{Embedding track.}
We extract the final-layer hidden state at every word, mean-pooled over its subword tokens (sentence datasets use pooled sentence features, and continuous recordings require temporal alignment, Appendix~\ref{app:implementation}). A ridge regression maps these embeddings to the human response, with the regularization strength chosen by nested cross-validation and all scores computed on held-out folds (stories, passages, sentences or time chunks, depending on the dataset; Table~\ref{tab:ridge-config}). The score is the held-out Pearson correlation between predicted and observed responses. The embedding track uses the full representation, so it measures how much of the signal is linearly decodable from the model. The ridge encoders include no covariates: unlike the surprisal track, they are not given word length or frequency, so lexical information that a model's representation carries can contribute to the score.

\paragraph{Surprisal track.}
Surprisal is the model's next-token negative log probability, summed over the subword tokens of a word. For each dataset we fit a linear mixed-effects model with dataset-specific nuisance predictors (word length and frequency for reading data, plus the covariates required by each neural recording design) and compare it with the model that adds the current word's surprisal and, where specified, the previous word's surprisal as a spillover term. Random-effect structures are dataset-specific (Appendix~\ref{app:lme-formulas}). The score is the improvement in held-out log likelihood per test observation,
\begin{equation}
\Delta\ell=\frac{\ell_{\mathrm{full}}^{\mathrm{test}}-\ell_{\mathrm{base}}^{\mathrm{test}}}{n_{\mathrm{test}}}.
\label{eq:dll}
\end{equation}
A model whose surprisal adds nothing beyond the nuisance predictors scores 0, regardless of its size; negative values indicate worse held-out fit than the baseline. The observation is defined by the dataset and need not be a single word in the neural analyses. The two predictor tracks consequently measure different properties and use different outcome scales.

\subsection{Composite scores and group analysis}
\label{sec:design-stats}

\paragraph{Normalization, BrainScore and BehavScore.}
Raw scores (correlations, $\Delta\ell$) are not comparable across datasets: a correlation of 0.3 may be close to the achievable maximum for a noisy single-trial recording and far from it for a heavily averaged fMRI response. For embeddings, each dataset-unit score is therefore divided by its split-half correlation noise ceiling \citep{vanbree2025noiseceiling}; for surprisal, it is divided by a split-half $\Delta\ell$ reference built from the human data \citep{timkey2026dissociation} (Appendix~\ref{app:normalization}). We call the median normalized score across neural units a model's \emph{BrainScore} and the median across behavioral units its \emph{BehavScore}; each is defined for both predictor tracks. These ratios are not fractions of explained variance, and split-half likelihood references can be exceeded (Appendix~\ref{app:normalization}).

\paragraph{Pareto fronts.}
A model is Pareto-optimal if no other evaluated model scores at least as highly on both BrainScore and BehavScore and strictly higher on one. Front membership describes estimated scores without uncertainty-based selection.

\paragraph{Separate behavioral and neural regressions.}
For each predictor track we regress the normalized score of each (model, dataset unit) pair on the two constraint indicators, the size of the model, and its language-modeling quality on that dataset's text, separately within each response domain $m\in\{\mathrm{behavioral},\mathrm{neural}\}$:
\begin{align}
Y_{ij}^{(m)}={}&\beta_0^{(m)}+\beta_M^{(m)}M_i+\beta_D^{(m)}D_i+\beta_{MD}^{(m)}M_iD_i+\beta_S^{(m)}\log_{10}P_i+\beta_Q^{(m)}\log_{10}\mathrm{PPL}_{ij} \nonumber\\
&+u_j+b_{Sj}\log_{10}P_i+b_{Qj}\log_{10}\mathrm{PPL}_{ij}+f_{\ell(i)}+v_i+c_{k(j)}+\epsilon_{ij}.
\label{eq:group}
\end{align}
Here $Y$ is the noise-ceiling-normalized direct-prediction score of model $i$ on unit $j$, $P$ is parameter count, and PPL is perplexity on the dataset text (both centered). Standard LLMs are the reference. $\beta_M$ and $\beta_D$ are the contrasts of each single-constraint group with Standard LLMs, the combined contrast of Constrained-BabyLM models is $\beta_M+\beta_D+\beta_{MD}$, and $\beta_S$ and $\beta_Q$ absorb scale: a constraint is associated with an advantage only if its contrast is positive \emph{beyond} what size and perplexity predict. An interaction measures departure from additivity, not the combined group's advantage over the reference.

The random effects model the dependence in the data. Scores of one dataset unit share a random intercept $u_j$ and unit-specific slopes $b_{Sj}$ and $b_{Qj}$ for size and perplexity, because the association of size with alignment varies in strength and sign across datasets (Appendix~\ref{app:rsa-panels}). Models of one lineage share a random intercept $f_{\ell(i)}$: a lineage is a model family (e.g.\ all ALiBi models, whatever their training data, or all Pythia models), except that the eleven BabyLM models, which come from five independent submissions, are split into those submissions; there are 25 lineages. Each model has its own random intercept $v_i$, since the same model contributes one score per dataset unit, and behavioral units recorded in one corpus (first-pass, go-past and scan-path times) share a corpus intercept $c_{k(j)}$. Models are fit by restricted maximum likelihood with \texttt{lme4}, and tests use Satterthwaite degrees of freedom. The analysis includes the 115 models with usable scores and perplexities (56 Standard LLMs, 24 Memory-Constrained, 11 BabyLM, and 24 Constrained-BabyLM; the three OLMo-2 models are excluded because their perplexities were not computed), giving 1,822 behavioral and 1,257 neural rows in each track.

\paragraph{Testing the behavioral--neural difference.}
Separate fits show whether a constraint is associated with each domain, but a significant coefficient in one domain and a nonsignificant one in the other is not evidence that they differ. We therefore also fit a joint model of both domains, in which every fixed effect of Equation~\ref{eq:group} interacts with domain, with the same random effects plus lineage-by-domain and model-by-domain intercepts. The domain interactions directly estimate the difference between neural and behavioral coefficients.

\paragraph{Robustness.}
We assess how much the conclusions depend on the analysis choices in five ways (Appendix~\ref{app:robustness}): (i) a specification with only a dataset-unit random intercept, which treats the rows as otherwise independent; (ii) a cluster bootstrap that resamples model lineages with replacement within each constraint group; (iii) refits that exclude the dataset units whose likelihood references are flagged as unstable (Appendix~\ref{app:ceiling-values}), winsorize scores within unit, exclude the GPT-BERT models, or add model depth as a covariate; (iv) leave-one-unit-out refits; and (v) the same models on a within-unit robust z-score, $(x-\mathrm{median})/(\mathrm{IQR}/1.349)$ over models, which expresses effects in units of the between-model spread and does not depend on the noise ceilings. We base our conclusions on effects that hold across these checks. The groups still differ in architecture, objective and training corpus beyond the two binary factors, so the estimates are adjusted associations, not causal effects.

\subsection{Representational evaluation}
\label{sec:design-rsa}

Ridge encoding models have a number of free parameters that grows with the hidden size of the model, so a score that increases with model size may partly reflect the capacity of the regression. We therefore complement direct prediction with two analyses of representational dissimilarity matrices (RDMs; \citealp{kriegeskorte2008rsa}), in which a model contributes a single stimulus-by-stimulus matrix whatever its width. Embedding distances are one minus the Pearson correlation between final-layer vectors after standardizing each feature across stimuli \citep{timkey2021allbark}. Surprisal distances are absolute differences in surprisal. Response distances are correlation distances for neural patterns and absolute differences for scalar reading-time or N400 responses.

Unfitted Pearson RSA correlates corresponding model and response distances without learning a response mapping. RDM regression instead measures the held-out Gaussian log-likelihood improvement when a model RDM is added to word-length and frequency RDMs, mirroring Equation~\ref{eq:dll}. Training and test partitions separate stimuli rather than randomly separating pairs \citep{schutt2023statistical}. Both analyses yield ceiling-normalized behavioral and neural composites, but their numerical scales differ from direct prediction. The representational fronts cover all 118 models, 16 behavioral measures, and eight neural datasets; Le Petit Prince, Natural Stories fMRI, and Fedorenko are absent. Appendix~\ref{app:rsa} specifies the scoring and aggregation.
